% !TEX program = lualatex
\ifdefined\mobileformat
  \documentclass[12pt]{extarticle}
\else
  \documentclass[20pt]{extarticle}
\fi
\ifdefined\mobileformat
  \usepackage[left=0.1in, right=0.1in, top=0.25in, bottom=0.35in, headsep=2pt]{geometry}
  \setlength{\headheight}{14pt}
\else
  \usepackage[left=0.1in, right=0.1in, top=0.5in, bottom=0.6in, headsep=4pt]{geometry}
  \setlength{\headheight}{20pt}
\fi
\usepackage{enumitem}
\usepackage{titlesec}
\ifdefined\mobileformat
  \setlist{nosep, leftmargin=1.2em}
  \titleformat{\section}{\normalsize\bfseries}{}{0em}{}
\else
  \setlist{nosep}
  \titleformat{\section}{\large\bfseries}{}{0em}{}
\fi
\titlespacing*{\section}{0pt}{1ex}{0.6ex}
\raggedright
\setlength{\parindent}{0pt}
\setlength{\parskip}{0.5ex plus 0.2ex}
\setcounter{secnumdepth}{0}

% --- fonts via fontspec (loads .ttf directly, preserving TT hinting) ---
\usepackage{fontspec}
\setmainfont{Gentium Book Plus}[
  BoldFont = {Gentium Book Plus Bold},
  ItalicFont = {Gentium Book Plus Italic},
  BoldItalicFont = {Gentium Book Plus Bold Italic},
]

\newcommand{\doctitle}{Yet Another Theory of Programming}
\usepackage{fancyhdr}
\pagestyle{fancy}
\fancyhf{}
\renewcommand{\headrulewidth}{0pt}
\fancyhead[L]{\scriptsize\doctitle}
\fancyhead[R]{\scriptsize\thepage}
\fancypagestyle{firstpage}{\fancyhf{}\renewcommand{\headrulewidth}{0pt}\fancyhead[R]{\scriptsize\thepage}}

\title{\doctitle}
\author{Jeremy Jacobson}
\date{March 9, 2026}

\begin{document}

\maketitle
\thispagestyle{firstpage}

\begin{quote}
\itshape
``It is important to have an appropriate understanding of what
programming is. If our understanding is inappropriate, we will
misunderstand the difficulties that arise in the activity, and our
attempts to overcome them will give rise to conflicts and
frustrations.''
\upshape
\hfill --- Peter Naur, \emph{Programming as Theory Building} (1985)
\end{quote}

\bigskip

\begin{abstract}
While I am certain that much has been written on topics adjacent to
those discussed here, I have chosen not to review the literature beyond the references cited
herein. I am writing this because I found that, in order to use AI
effectively for coding, I needed such a thing, a theory of
programming.
\end{abstract}

\section{Introduction}

In the essays collected in \emph{The Mythical Man-Month}
\cite{brooks1995}, Frederick Brooks chooses to decompose the practice
of programming into two parts, which, following Aristotle, he calls
the \emph{essence} and the \emph{accident}. The essence has as a core
element what Brooks calls \emph{conceptual integrity}, the design of
the thing, the unity of its ideas. The accident encompasses all of
the typical activities we know as programming: in particular, an
implementation of the concept in some language, on some machine, under
some set of constraints.

Elsewhere in that book, Brooks recalls a notion due to Dorothy Sayers
\cite{sayers1941} regarding the creative process. Sayers holds that
creation, no matter which domain we are discussing, whether art,
science, programming, or anything else, is made up of three
fundamental parts:
\begin{enumerate}
  \item The \emph{theorization}: the idea, conceived in the mind
    before any material work begins.
  \item The \emph{implementation}: the putting of pen to paper, so to
    speak: writing the code, putting the brush to the canvas.
  \item The \emph{use}: the produced object being received and used by
    others.
\end{enumerate}

It is only when the work is used, whether it is art that is
being viewed or a program that is being run, that the concept is
really understood.

All of this was captured nicely by Peter Naur, known for the
Backus--Naur Form (BNF), in his essay \emph{Programming as Theory
  Building}, written in 1985 \cite{naur1985}. For me, the key quote
is this:
\begin{quote}
  ``On the theory-building view, the primary result of the programming
  activity is the theory held by the programmers. Since this theory,
  by its very nature, is part of the mental possession of each
  programmer, it follows that the notion of the programmer as an
  easily replaceable component in the program production activity has
  to be abandoned.''
\end{quote}

Later, Naur says in
his conclusion:
\begin{quote}
  ``Further, on this view, the notion of a programming method
  understood as a set of rules of procedure to be followed by the
  programmer is based on invalid assumptions and so has to be
  rejected.''
\end{quote}

\section{The Limits of Specifications and Standards}
One might hope, and many currently seem to believe, that with sufficient context, specifications and procedural standards, AI will be able to contribute to
software development at a level that meets or exceeds the most skilled developers.

However, if one takes the theory building view of programming, then clearly there remain challenges with such an approach. In the words of Naur:

\begin{quote}
  ``For a new programmer to come to possess an existing theory of a
  program, it is insufficient that he or she has the opportunity to
  become familiar with the program text and other documentation. What
  is required is that the new programmer has the opportunity to work
  in close contact with the programmers who already possess the
  theory, so as to be able to become familiar with the place of the
  program in the wider context of the relevant real-world situations,
  and so as to acquire the knowledge of how the program works and how
  unusual program reactions and program modifications are handled
  within the program theory.''
\end{quote}

Later, Naur writes:
\begin{quote}
  ``A very important consequence of the theory-building view is that
  program `revival', i.e., re-establishing the theory of a program
  merely from the documentation, is strictly impossible.''
\end{quote}

Indeed, consider the example of research mathematics. Despite no shortage of articles and textbooks, mathematics at the highest levels is taught only in a student advisor relationship. To reach an understanding of the 'theory of the program' when it comes to research mathematics, with few exceptions it is necessary for a student to work for 4-6 years under a working mathematician. 

\section{Capturing the Theory}

Recall that program maintenance may be broken into three types:
\begin{enumerate}
  \item \textbf{Corrective}, which involves fixing bugs.
  \item \textbf{Perfective}, which pertains to program optimizations.
  \item \textbf{Adaptive}, which refers to program modifications
    demanded by changing external circumstances.
\end{enumerate}

For AI to be a full participant in the software development lifecycle, it must be able to perform the adaptive maintenance that is often the most expensive drain on a developer's time. 

As Naur argues, and we agree, even the most elaborate set of standards
and specifications for the software development process will fail to
capture the theory of programming. In turn, in our use of AI, this
will lead to AI developers (that is, the AI itself) being unable to
perform specifically the adaptive maintenance that is so critical and
truly separates software engineering from other engineering
disciplines as it requires that software behave more like a living
thing, able to adapt to a changing environment.

\section{Actions to Be Taken}

We take this to mean that if one starts from a well-accepted
methodology, for instance, the IEEE standards or, in our case, our
personal favorite, the MIL--STD--498 standards and
specifications \cite{milstd498}, the latter consisting of proven templates
that have been used for decades and are still employed today to
engineer the software development effort on countless systems
including large and critical ones, then one finds that a key document
is the so-called \emph{Operational Concept Description} (OCD).

This document describes the general nature of the system upon which
development is to be done. It summarizes the history of system
development, operation, and maintenance; identifies the project
sponsor, acquirer, user, developers, and support organizations;
identifies environments and sites of operation. It includes justification for and the nature of proposed changes. This even includes a
subsection called Changes Considered but Not Included. The
justification asks for some provision of deficiencies or limitations
in the current system or situation that prevent an adequate response to
these factors.

For us, the authors of the operating concept were certainly well aware
of the issues that Naur raised. I believe they would have, if it were
possible at the time, implemented what he is describing. He
essentially asserts that in order for someone to program on a code base that is new to them, they first need a programmer who built the
original system to teach them about it, akin to how a musician teaches
a student music or a writer teaches someone how to write. This
includes discussions of the relation between the program and the
relevant aspects and activities of the real world, and of the limits
set on the real-world matters dealt with by the program. Indeed, these
are the words of Naur.

Of course, no such thing was possible at the time that these military
standards were written. In fact, no such thing has been possible until
today.

AI now has the ability to live in a code base from the first iteration. What we are
suggesting is that we go from an operating concept \emph{document} to an operating concept \emph{agent}: a professor of
the theory of programming for that particular code base.

Further, we have at our disposal important elements that were not available in
the past. Previously, we had the code itself and written
documentation. Now, an agent can review the Git history. It can
review all of the pull requests, all the issues that were raised, the comments; it
can analyze how the system has actually responded to each nuance of
situation based on the observed behavior. The lead engineer on the code base, the person who was responsible for conceptual
integrity, their actions, commits, comments on pull requests, etc., all of this taken together may allow us to recover to a sufficient level of fidelity the theory of the program for that code base.

We feel this is the final keystone in the building of an AI
programming system for software engineering.

\begingroup
\setlength{\parskip}{0pt}
\begin{thebibliography}{9}

\bibitem{brooks1995}
Frederick P.\ Brooks, Jr., \emph{The Mythical Man-Month: Essays on
  Software Engineering}. Anniversary Edition. Addison-Wesley, 1995.

\bibitem{milstd498}
U.S.\ Department of Defense, \emph{MIL--STD--498: Software
  Development and Documentation}. 1994.

\bibitem{naur1985}
Peter Naur, ``Programming as Theory Building.''
\emph{Microprocessing and Microprogramming}, 15(5):253--261, 1985.

\bibitem{sayers1941}
Dorothy L.\ Sayers, \emph{The Mind of the Maker}. Methuen, 1941.

\end{thebibliography}
\endgroup

\end{document}
